% Einheit 3 von 5 – Kumulierte Wahrscheinlichkeiten (bank/binomialverteilung/e3.jsonl, zusammenbau v0.1)
%% TODO Zeitmarke Einheit 3: Marken-Zeile nicht lesbar
%% TODO Zweigzeile Teil 1: was hier gelernt wird (Satz in Schülersprache aus dem Einheitstitel) – Einheit 3
\einheitenkopf[e3]{Einheit 3 von 5 · Kumulierte Wahrscheinlichkeiten}
\zweigzeile{Abitur GK}
\setcounter{aufgabe}{14}

%% TODO Titel: Ich-kann-Satz fehlt in der Bank (Platzhalter: Kettenname) – Nr. 15
%% TODO Anweisung: ein Satz über der ersten Teilaufgabe fehlt in der Bank – Nr. 15
\begin{aufgabe}{Kumulierte Wahrscheinlichkeiten}
%% TODO \rechenplatz{n}: Schrittzahl des Grundfalls fehlt in der Bank (nur bei fester Schreibform, 2.3 b)
\begin{teile}
\teil Ein Würfel wird 20-mal geworfen; X ist die Anzahl der Sechsen. „Weniger als fünf Sechsen“ – welcher Ansatz passt? \kreuz{$P(X \le 4)$} \\ \kreuz{$P(X \le 5)$} \\ \kreuz{$1 - P(X \le 4)$}
\teil 30 Würfe; X ist die Anzahl der Treffer. Übersetze „höchstens 6 Treffer“ in einen Ansatz mit $P(X \le \ldots)$.
\teil 35\,\% der Schüler einer Schule kommen mit dem Bus. 25 Schüler werden zufällig ausgewählt. Höchstens 7 Busfahrer – mit welcher Wahrscheinlichkeit? Runde auf vier Stellen. \leerfeld
\teil 45\,\% der Kunden eines Cafés bestellen Kuchen. An einem Nachmittag kommen 30 Kunden. Mindestens 15 bestellen Kuchen – mit welcher Wahrscheinlichkeit? Runde auf vier Stellen. \leerfeld
\teil Im Mittel ist eine von acht Glühbirnen einer Sorte defekt; die Anzahl X der defekten unter zufällig ausgewählten gilt als binomialverteilt. 60 Birnen werden ausgewählt. A: höchstens 6 sind defekt. B: mehr als 8 und weniger als 12 sind defekt. Berechne P(A) und P(B) \hfill (Abitur 2018 GK).
\teil 7\,\% der Eier eines Hofes sind Knickeier; die Anzahl X der Knickeier gilt als binomialverteilt. 60 Eier werden zufällig ausgewählt. A: genau drei sind Knickeier. B: mindestens 10\,\% der Eier sind Knickeier. Berechne P(A) und P(B) \hfill (Abitur 2018 LK).
\teil 35\,\% der Mitarbeiter einer Firma arbeiten im Homeoffice, die übrigen im Büro. 60 Mitarbeiter werden zufällig ausgewählt; H ist die Anzahl im Homeoffice. Mit welcher Wahrscheinlichkeit ist die Anzahl im Büro mindestens doppelt so groß wie die im Homeoffice? \hfill (Abitur 2023 GK) \\ \leerfeld
\teil 12\,\% der Erwachsenen eines Landes haben Heuschnupfen. Für eine Studie werden 1\,500 Erwachsene zufällig ausgewählt; X ist die Anzahl mit Heuschnupfen. Bestimme in einem geeigneten Modell die Wahrscheinlichkeit, dass X um mehr als 10\,\% vom Erwartungswert abweicht \hfill (Abitur 2017 LK). \leerfeld
\teil 6\,\% der Fahrräder, die die Polizei kontrolliert, haben kein funktionierendes Licht. 80 Fahrräder werden kontrolliert. Gib die Bedeutung des Terms $1 - (0{,}94^{80} + 80 \cdot 0{,}06 \cdot 0{,}94^{79})$ im Sachzusammenhang an \hfill (Abitur 2017 LK).
\teil Ein Glücksrad zeigt Stern mit der Wahrscheinlichkeit $0{,}6$ und Kreis mit $0{,}4$; es wird achtmal gedreht und die Folge der acht Symbole notiert. Gewinn gibt es bei mehr als vier Kreisen, einen Extrapreis bei genau zwei Sternen unter den vorderen sechs Symbolen. Berechne die Wahrscheinlichkeit, Gewinn und Extrapreis zu erhalten \hfill (Abitur 2022 LK). \leerfeld
\end{teile}
\end{aufgabe}

%% TODO Titel: Ich-kann-Satz fehlt in der Bank (Platzhalter: Kettenname) – Nr. 16
%% TODO Anweisung: ein Satz über der ersten Teilaufgabe fehlt in der Bank – Nr. 16
\begin{aufgabe}{Bedingte Restwahrscheinlichkeit nach bekannten Ergebnissen über die Binomialverteilung berechnen}
\begin{teile}
\teil Eine Münze wird 400-mal geworfen. Nach den ersten 100 Würfen ist 58-mal Zahl gefallen. Mit welcher Wahrscheinlichkeit fällt bei höchstens 52\,\% aller 400 Würfe Zahl? \hfill (Abitur 2023 GK) \\ \leerfeld
\end{teile}
\end{aufgabe}

%% TODO Titel: Ich-kann-Satz fehlt in der Bank (Platzhalter: Kettenname) – Nr. 17
%% TODO Anweisung: ein Satz über der ersten Teilaufgabe fehlt in der Bank – Nr. 17
\begin{aufgabe}{Wahrscheinlichkeit für zwei unabhängige Spieler als Produkt binomialer Wahrscheinlichkeiten berechnen}
\begin{teile}
\teil Anna trifft beim Freiwurf mit der Wahrscheinlichkeit $0{,}6$, Ben mit $0{,}5$; beide werfen 10-mal. Mit welcher Wahrscheinlichkeit haben beide mindestens 6 Treffer? Runde auf vier Stellen. \leerfeld
\end{teile}
\end{aufgabe}

%% TODO Titel: Ich-kann-Satz fehlt in der Bank (Platzhalter: Kettenname) – Nr. 18
%% TODO Anweisung: ein Satz über der ersten Teilaufgabe fehlt in der Bank – Nr. 18
\begin{aufgabe}{Kumulierte Wahrscheinlichkeiten – Fehler finden, Begründen, Anwendung, Darstellungswechsel}
\begin{teile}
\teil Ein Glücksrad bringt mit der Wahrscheinlichkeit $0{,}35$ einen Gewinn und wird 20-mal gedreht; X ist die Anzahl der Gewinne. Paul rechnet für „mehr als 7 Gewinne“: \rechnung{P &= 1 - P(X \le 6)} Finde den Fehler.
\teil Begründe, warum man „mindestens k Treffer“ als eins minus „höchstens k − 1 Treffer“ berechnet und nicht als eins minus „höchstens k Treffer“.
\teil Ein Hotel hat 120 Zimmer und nimmt 125 Buchungen an, weil 6\,\% der Buchungen kurzfristig storniert werden. Mit welcher Wahrscheinlichkeit reichen die Zimmer? Runde auf vier Stellen.
\teil 15\,\% der Pflanzen einer Sorte blühen weiß; 25 Pflanzen werden gesetzt. Schreibe die Wahrscheinlichkeit für „höchstens 3 blühen weiß“ mit dem Summenzeichen und berechne sie auf vier Stellen.
\end{teile}
\end{aufgabe}
